\documentclass[article,screen,12pt, nonacm]{acmart}
\let\Bbbk\relax %
\usepackage{standalone}
\usepackage{vmlmacros}
\usepackage{titlesec}
\titlespacing*{\section}{0pt}{0.7em plus 0.2em minus 0.1em}{0.5em}
\usepackage{titlesec}

% Reduce whitespace before and after section titles
\titlespacing*{\section}
  {0pt}    % left margin
  {0.6em plus 0.2em minus 0.1em} % before section
  {0.4em plus 0.1em minus 0.1em} % after section

\titlespacing*{\subsection}
  {0pt}
  {0.5em plus 0.1em minus 0.1em}
  {0.3em plus 0.1em minus 0.1em}

\titlespacing*{\subsubsection}
  {0pt}
  {0.4em plus 0.1em minus 0.1em}
  {0.2em plus 0.05em minus 0.05em}

% Slightly reduce paragraph spacing to match
\setlength{\parskip}{0pt}
\setlength{\parindent}{1em}
\AtBeginDocument{%
  \providecommand\BibTeX{{%
    \normalfont B\kern-0.5em{\scshape i\kern-0.25em b}\kern-0.8em\TeX}}}
\usepackage{outlines}
\setlength{\headheight}{14.0pt}
\setlength{\footskip}{13.3pt}
\citestyle{acmauthoryear}
\begin{document}
\title{An Alternative to Pattern Matching, Inspired by Verse}
\author{Roger Burtonpatel}
\email{roger.burtonpatel@tufts.edu}
\affiliation{%
  \institution{Tufts University}
  \streetaddress{419 Boston Ave}
  \city{Medford}
  \state{Massachusetts}
  \country{USA}
  \postcode{02155}
}
\bibliographystyle{plainnat}




\renewcommand{\shortauthors}{Burtonpatel}

\rab{fix odd page indentation}
\begin{abstract}
    Pattern matching %
    appeals %
    to functional programmers for its expressiveness and good cost model, 
    but 
    it %
    fails to express %
    certain computations without verbosity. Verbosity %
    can be reduced %
    by using extended forms of pattern matching,
    but extensions %
    are %
    not standardized, and making certain extensions like or-patterns and pattern
    guards cooperate in a large-scale programming language is a significant
    implementation challenge that has not seen success until recent years. 
    
    
    Alternatively to pattern matching, computations can be expressed succinctly
    by using equations, as has been demonstrated by the Verse programming
    language. But such computations may have time and space costs that can be
    hard to predict. 
    As a compromise, 
    we %
    ~propose %
    a new language,~\VMinus, %
    which~uses equations in a limited way that makes time and space costs easy
    to predict. 
    In
    comparative examples,
    \VMinus~ %
    expresses %
    computations as succinctly as pattern matching with popular extensions, and just like pattern matching, it %
    can be compiled %
    to a decision tree. 
    
    To evaluate %
    the succinctness of \VMinus,
    we compare %
    it with a novel pattern-matching language, \PPlus, which exhibits
    three popular pattern matching extensions 
    in unison: side conditions, or-patterns, and pattern guards. 
    We evaluate %
    the expressiveness of \VMinus using four smaller contrived examples 
    and a larger example from \_. 

    An implementation of~\VMinus, \PPlus, and the match compiler can be found
    at~\url{https://github.com/rogerburtonpatel/vml}.
\end{abstract}

\maketitle

\section{Introduction}

Perhaps the most beloved tool among functional programmers for examining and
deconstructing data is pattern matching. 
Pattern matching is also an established and well-researched topic
\citep{wadler1987views, macqueen1985tree, burton1993pattern, palao1996new,
maranget, bpc}. It is appreciated by programmers and researchers alike for two
main reasons: It enables~\it{implicit} data deconstruction, and it has a
desirable cost model. Specifically (regarding the latter), pattern matching can
be compiled to a~\it{decision tree}, a data structure that enforces linear
runtime performance by guaranteeing no part of the data will be examined more
than once.~\citep{maranget}

However, pattern matching cannot express certain common computations succinctly,
forcing programmers who wish to express these computations to duplicate code,
nest~\it{case} expressions, and create multiple points of truth. To mitigate
this, designers of popular programming languages have introduced~\it{extensions}
to pattern matching, most commonly \textit{side conditions}(cite),
\textit{or-patterns}(cite), and \textit{pattern guards}(cite).
\rab{i define these below\dots do so here or in the main paper?}

Extensions strengthen pattern matching, but they are not standardized, so each
popular programming language with pattern matching features its own unique suite
of extensions. Rather than continuing to extend pattern
matching~\it{ad hoc}, a worthwhile goal could be to find an alternative that
doesn't need extensions. A tempting possibility was introduced by the
programming language Verse~\citep{verse}. In Verse, a programmer can deconstruct
data using ~\it{equations}. Equations are flexible, and it appears that they can express everything that
pattern matching can, including the popular extensions. 

But a full implementation of Verse is complicated, cost-wise. Verse is a
functional logic programming language, and expressions can backtrack at runtime
and return multiple results, both of which are hard to predict in their costs. 


In this paper, we~show that the expressive quality of Verse's equations and the
decision-tree property of patterns can be combined in a single language. Since
the language is a streamlined adaptation of Verse with a reduced feature set,
we~call it~\VMinus (“V minus”). For comparison, we employ~\PPlus, a novel
pattern-matching language with three popular extensions. We compare~\VMinus
with~\PPlus in several small examples and a large example which is \padlexample. 


\rab{todo flow}

Our contributions are: 
We~have formalized~\PPlus and \VMinus with a big-step operational semantics
(Section~\ref{vminus}), we~have formalized decision trees into a core
language~\D (“D”) with a big-step operational semantics (Section~\ref{d}), 
we~have formalized a translation from~\VMinus to~\D (Sections~\ref{vminustod}),
and we~have implemented both languages in Standard ML. 



\section{Pattern Matching and its Extensions}
\label{pmandextensions}
    \textbf{\rab{GOAL: Introduce pattern matching in a way that's neither
    dismissive of the interesting problem nor insulting to the audience's
    intelligence and knowledge of functional programming.}}


    Pattern matching is a battle-tested and beloved features in declarative and
    functional programming, but it often requires extensions to simplify cases
    that are otherwise troublesome. 
    
    
    Specifically, in sophisticated programs that involve complex
    decision-making, pattern matching without extensions tends to duplicate
    code. Consider this example of "simple" pattern matching (without
    extensions): 

    \rab{Example; will discuss which to use. Adaptation of 'cellar' example
    seems good}.

    \rab{Example says along the lines of: "Without extensions to pattern
    matching, expressing the algorithm looks like this: \dots}. 

    The code gets the job done, but it is cumbersome to write and read because
    of the large amount of duplication, which arises from the limitations of
    what can be expressed in simple pattern matching. 
    
    Furthermore, if the spec changes, will the programmer be shrewd enough to
    find and modify all the duplicated parts of the program to match the change?
    Such an example demands an improvement. Indeed, we can do better, by 
    augmenting pattern matching with \textit{extensions}. 


    \subsection{Popular extensions to pattern matching}
    \label{extensions}




% ------------------------------------------------------------



% ------------------------------------------------------------
    


    \rab{TODO: Format this paragraph so it presents information as clearly as possible.}
    Duplicated code, especially in a complex case expression, can make code
    difficult to maintain and more bug-prone. To address this, many large-scale
    functional programming languages have added extensions to pattern matching,
    most commonly \textit{side conditions}(cite), which extend a pattern match with a
    boolean guard, \textit{or-patterns}(cite), which allow multiple branches in a case
    expression to share a right-hand side, and \textit{pattern guards}(cite), which
    allow multiple sequential pattern-matches and computations to guard a
    right-hand side. 

    \rab{Side condition example}
    Side conditions do \dots which eliminate the need to \dots, 
    illustrated in this small example of measuring shapes: 

        \rab{The examples are funny, but they should be more serious for this 
        submission.}

        \rab{TODO: merge figures, make them appear in the right place}
      
          \begin{figure}[h]
\centering
\begin{minipage}{0.45\linewidth}
\small
\begin{verbatim}

    type standard_shape = Square of float 
                    | Triangle of float * float
                    | Trapezoid of float * float * float

let exclaimTall sh =
  match sh with 
    | Square s -> 
        if s > 100.0 
        then "Wow! That's a sizeable square!"
        else "Your shape is small." 
    | Triangle (_, h) -> 
        if h > 100.0 
        then "Goodness! Towering triangle!"
        else "Your shape is small." 
    | Trapezoid (_, _, h) -> 
        if h > 100.0
        then "Zounds! Tremendous trapezoid!"
        else "Your shape is small." 
\end{verbatim}
\caption*{Without extensions}
\end{minipage}
\hfill
\begin{minipage}{0.45\linewidth}
\small
\begin{verbatim}
let exclaimTall sh =
  match sh with 
    | Square s when s > 100.0 ->
        "Wow! That's a sizeable square!"
    | Triangle (_, h) when h > 100.0 ->
        "Goodness! Towering triangle!"
    | Trapezoid (_, _, h) when h > 100.0 -> 
        "Zounds! Tremendous trapezoid!"
    | _ -> "Your shape is small." 
\end{verbatim}
\caption*{With side conditions}
\end{minipage}
\caption{Comparing pattern matching with and without side conditions.}
\end{figure}


\begin{figure}[ht]
\centering
\begin{minipage}{0.47\linewidth}
\small
\begin{verbatim}
  let tripleLookup (rho : finitemap) (x : int) =
      match lookup rho x with Some w -> 
        (match lookup rho w with Some y -> 
          (match lookup rho y with Some z -> z
          | _ -> handleFailure x)
        | _ -> handleFailure x)
      | _ -> handleFailure x
\end{verbatim}
\caption*{OCaml version with nested pattern matching}
\end{minipage}
\hfill
\begin{minipage}{0.47\linewidth}
\small
\begin{verbatim}
  tripleLookup rho x
     | Just w <- lookup rho x
     , Just y <- lookup rho w
     , Just z <- lookup rho y
     = z
  tripleLookup _ x = handleFailure x
\end{verbatim}
\caption*{Haskell version with pattern guards}
\end{minipage}
\caption{Comparing \texttt{tripleLookup} using pattern guards.}
\label{fig:pmtriplelookup}
\Description{Side-by-side comparison of tripleLookup implementations in OCaml and Haskell.}
\end{figure}
\label{fig:guardtriplelookup}

    \rab{Or pattern example}
    Or patterns do \dots which eliminate the need to \dots

     \begin{figure}[h]
\centering
\begin{minipage}{0.45\linewidth}
\small
\begin{verbatim}
let describeColor c =
  match c with
    | Red -> "Warm color"
    | Orange -> "Warm color"
    | Yellow -> "Warm color"
    | Blue -> "Cool color"
    | Green -> "Cool color"
\end{verbatim}
\caption*{Without extensions}
\end{minipage}
\hfill
\begin{minipage}{0.45\linewidth}
\small
\begin{verbatim}
let describeColor c =
  match c with
    | Red | Orange | Yellow -> "Warm color"
    | Blue | Green -> "Cool color"
\end{verbatim}
\caption*{With or-patterns}
\end{minipage}
\caption{Comparing pattern matching with and without or-patterns.}
\end{figure}

    \newpage
    Individually, each of these three extensions grants a powerful abstraction
    that reduces code duplication. It follows that they should retain their
    powers of brevity when operating in unison, and indeed this is the case:
    with these three extensions co-operating, the previous example reduces like
    this \dots

    Given the expressive power of this combination and each extension's
    individual popularity, it would seem obvious that they are all supported by
    myriad functional programming languages. But that's not true: among Haskell,
    OCaml, Scala, Erlang/Elixir, Rust, F\#, Rocq, and Agda~\citep{haskell,
    ocaml, scala, erlang, elixir, rust, fsharp, agda}, only Haskell has all
    three of side conditions, or patterns, and pattern guards, and then only as
    of 2023! 

    - No strict language with all three + significant implementation effort in
    Haskell to get them all working -> evidence enough that having all of them
    is a non-trivial problem. 
    
    We~find the extension story somewhat unsatisfying. At the very least,
    we~want to be able to use pattern matching, with the extensions we~want, in
    a single language. Or, we~want an alternative that gives me the expressive
    power of pattern matching with these extensions. 

    \section{Equations}

    An intriguing alternative to pattern matching exists in~\it{equations}, from
    the Verse Calculus (\VC), a core calculus for the functional logic
    programming language~\it{Verse}~\citep{antoy2010functional,
    hanus2013functional, verse}. (For the remainder of this work, we~use “\VC”
    and “Verse” interchangeably.)

    Verse's equations and choice mechanism allow for much of the same expressive
    power as pattern matching with extensions. Rather than go into too much
    detail (the curious may examine the Verse paper~\citep{verse}), we support
    this claim with a familiar example: 

    \rab{TODO : insert existing examples}

    The brevity of the code in the figures is is promising for~\VC. If it rivals
    pattern matching with popular extensions in expressiveness, and~\VC
    does everything using only equations and choice, it seems like the language
    is an intriguing option for writing code! 


    But the VC semantics has a cost model that is hard to predict. Why? In~\VC,
    names (logical variables) are~\it{values}, and they can just as easily be
    the result of evaluating an expression as an integer or tuple. To bind these
    names,~\VC, like other functional logic languages, relies
    on~\it{unification} of its logical variables and \it{search} at runtime to
    meet a set of program constraints~\citep{antoy2010functional,
    hanus2013functional}. Combining unifying logical variables with search
    requires backtracking, which can lead to exponential runtime
    cost~\citep{hanus2013functional, wadler1985replace, clark1982introduction}. 


    Pattern matching, by contrast, can be compiled to a~\it{decision tree}, a
    data structure that enforces linear runtime performance by guaranteeing no
    part of the scrutinee will be examined more than once~\citep{maranget}. A
    decision tree does not backtrack: once a program makes a decision based on
    the form of a value, it does not re-test it later with new information. 
    

    Figures: 

    \begin{figure}[ht] 
        \begin{minipage}[h]{0.54\linewidth}
          \verselst
          \begin{lstlisting}[numbers=none, basicstyle=\tiny, xleftmargin=.2em,
                            showstringspaces=false,
                            frame=single]
$\exists$ exclaimTall. exclaimTall = $\lambda$ sh. 
  one { 
      $\exists$ s. sh = $\langle$Square s$\rangle$; 
      s > 100.0; "Wow! That's a sizeable square!"
    | $\exists$ w h. sh = $\langle$Triangle, w, h$\rangle$; 
      h > 100.0; "Goodness! Towering triangle!"
    | $\exists$ b1 b2 h. sh = $\langle$Trapezoid, b1, b2, h$\rangle$;
      h > 100.0; "Zounds! Tremendous trapezoid!"
    | "Your shape is small." }
            \end{lstlisting}
                \Description{exclaimTall}
        \subcaption{\tt{exclaimTall} in~\VC }
            \label{fig:verseexclaimtall} 
        \end{minipage}%
        \begin{minipage}[h]{0.5\linewidth}
          \verselst
          \begin{lstlisting}[numbers=none, basicstyle=\tiny, xleftmargin=2em,
                        frame=single]
$\exists$ tripleLookup. tripleLookup = $\lambda$ rho x. 
    one { $\exists$ w. lookup rho x = $\langle$Just w$\rangle$; 
          $\exists$ y. lookup rho w = $\langle$Just y$\rangle$; 
          $\exists$ z. lookup rho y = $\langle$Just z$\rangle$;
          z
        | handleFailure x }
          \end{lstlisting}
                \Description{exclaimTall}
            \subcaption{\tt{tripleLookup} in~\VC }
            \label{fig:versetriplelookup} 
        \vspace{4ex}
        \end{minipage} 
        \begin{minipage}[h]{\linewidth}
          \verselst
          \begin{lstlisting}[numbers=none, basicstyle=\tiny, xleftmargin=9em, 
                            frame=single, showstringspaces=false]
$\exists$ game_of_token. game_of_token = $\lambda$ token. 
  $\exists$ f. one {
         token = one { $\langle$BattlePass, f$\rangle$ | $\langle$ChugJug, f$\rangle$ | $\langle$TomatoTown, f$\rangle$}; 
           $\langle$"Fortnite", f$\rangle$
       | token = one { $\langle$HuntersMark, f$\rangle$ | $\langle$SawCleaver, f$\rangle$}; 
           $\langle$"Bloodborne", 2 * f$\rangle$
       | token = one { $\langle$MoghLordOfBlood, f$\rangle$ | $\langle$PreatorRykard, f$\rangle$}; 
           $\langle$"Elden Ring", 3 * f$\rangle$
       |   $\langle$"Irrelevant", 0$\rangle$ }
          \end{lstlisting}
                \Description{game_of_token in~\VC}
        \subcaption{\tt{game\_of\_token} in~\VC }
            \label{fig:versegot} 
        \vspace{4ex}
        \end{minipage}%
    \caption{Code for the~\ref{pmandextensions} functions with equations looks
    similar, and it doesn't need extensions.}
    \label{fig:verseextfuncs}
      \end{figure}
        


      \rab{todo: name this better}
    \section{Two Languages, Two Solutions}
    We have created and implemented a semantics for three core languages that
    aim to bridge the gap between pattern matching, equations, and decision
    trees. They are \PPlus, which has pattern matching with all the above
    extensions, \VMinus, which has equations without backtracking or multiple
    results, and \D, which has decision trees. \PPlus\ and \VMinus serve to
    explore the design space of pattern matching and equations, while \D\ serves
    as the target of translation to reason about efficiency. Each language is an
    extension of the untyped lambda calculus with value constructors. 
    
    Like in \VC, every lambda calculus term is valid in our languages and has
    the same semantics. Also like the lambda calculus and \VC, all three
    languages are \it{strict}, meaning expressions are always evaluated when
    they are bound to variables. 
    

\subsection{Introducing \PPlus}
\label{pplus}

\begin{figure}
\begin{center}
\begin{bnf}
$P$ : \textsf{Programs} ::=
$\bracketed{d}$ : definition
;;
$d$ : \textsf{Definitions} ::=
| $\tt{val}\; x\; \tt{=}\; \expr$ : bind name to expression
;;
$\expr$ : Expressions ::= 
| $v$ : literal values 
| $x, y, z$ : names
| $K\bracketed{\expr}$ : value constructor application 
| $\ttbackslash x.\; \expr$ : lambda declaration  
| $\expr[1]\; \expr[2]$ : function application 
| $\tt{case}\; \expr\; \bracketed{p\; \ttrightarrow\; \expr}$ : case expression 
| \ttbraced{$\expr$}
;;
$p$ : \textsf{Patterns} ::= $p_{1}\pbar p_{2}$ : or-pattern
| $p \tt{,} p'$ : pattern guard 
| $p\; \tt{<-}\; \expr$ : pattern from explicit expression  
| $x$ : name 
| $\tt{\_}$ : wildcard 
| $K\; \bracketed{p}$ : value constructor application 
| $\tt{when}\; \expr$
| \ttbraced{$p$}
;;
$\v$ : Values ::= $K\bracketed{\v}$ : value constructor application 
| $\ttbackslash x.\; \expr$ : lambda value 
;;
$K$ : \textsf{Value Constructors} ::=
| $\tt{true}\; \vert\; \tt{false}$ : booleans
| $\tt{A-Z}x$ : name beginning with capital letter
% | $[\tt{-}\vert\tt{+}](\tt{0}-\tt{9})+$ : signed integer literal 
\end{bnf}
\end{center}
\Description{The concrete syntax of PPlus.}
\caption{\PPlus: Concrete syntax}
\label{fig:ppsyntax}
\end{figure}


    \PPlus\ packages common extensions to pattern matching with a standard core.
    In addition to bare pattern matching— names and applications of value
    constructors— the language includes pattern guards, or-patterns, and side
    conditions. 
    
    % Note that pattern guards subsume side conditions; we include side conditions
    % and separate them from guards for purpose of study. 
    

\subsection{Formal Semantics of \PPlus}














    
      

    
  Figure~\ref{fig:ppfuncs} provides an example of how a programmer might utilize
  \PPlus\ to solve the previous problems: 

    \begin{figure}[ht] 
      \begin{minipage}[h]{0.54\linewidth}
        \pplst 
        \begin{lstlisting}[numbers=none, basicstyle=\tiny, xleftmargin=.2em,
          showstringspaces=false,
          frame=single]
val exclaimTall = \sh.
  case sh of Square s, when (> s) 100 -> 
            Wow! That's A Tall Square!  
    | Triangle w h, when (> s) 100 ->
            Goodness! Towering Triangle!
    | Trapezoid b1 b2 h, when (> s) 100 -> 
            Zounds! Tremendous Trapezoid!
    | _ ->  Your Shape Is Small
  \end{lstlisting}
      \Description{exclaimTall in PPlus}
      \subcaption{\tt{exclaimTall} in \PPlus}
          \label{fig:ppexclaimtall} 
      \end{minipage}%
      \begin{minipage}[h]{0.5\linewidth}
        \pplst 
        \begin{lstlisting}[numbers=none, basicstyle=\tiny, xleftmargin=2em,
                      frame=single]
  val tripleLookup = \ rho. \x.
    case x of 
      Some w <- (lookup rho) x
    , Some y <- (lookup rho) w
    , Some z <- (lookup rho) y -> z
    | _ -> handleFailure x

   \end{lstlisting}
        \Description{tripleLookup in PPlus}
        \subcaption{\tt{tripleLookup} in \PPlus}
            \label{fig:pptriplelookup} 
        \vspace{4ex}
      \end{minipage} 
      \begin{minipage}[h]{\linewidth}
        \pplst 
        \begin{lstlisting}[numbers=none, basicstyle=\tiny, xleftmargin=9em,
          showstringspaces=false,
          frame=single]
val game_of_token = \t. 
  case t of  
    BattlePass f  | (ChugJug f | TomatoTown f) -> P (Fortnite f)
  | HuntersMark f | SawCleaver f               -> P (Bloodborne ((* 2) f))
  | MoghLordOfBlood f | PreatorRykard f        -> P (EldenRing ((* 3) f))
  | _                                          -> P (Irrelevant 0)
\end{lstlisting}
      \Description{game_of_token in PPlus}
      \subcaption{\tt{game\_of\_token} in \PPlus}
          \label{fig:ppgot}
      \vspace{4ex}
      \end{minipage}%
      \caption{The functions in \PPlus\ have the same desirable implementation
      as before.}
  \label{fig:ppfuncs}
    \end{figure}        
    
    

    
    


\subsection{Introducing \VMinus\ }
\label{vminus}

        To fuel the pursuit of smarter decision-making, we now draw inspiration
        from \VC. We begin by trimming away the aforementioned culprits that
        tends to lead to unpredictable or costly outcomes during runtime:
        backtracking and multiple results. Removing these language traits strips
        much of the identity of \VC; however, we do so in order to study only
        \VC's \it{equations} while staying grounded in an otherwise-typical
        programming context of single results and no backtracking at runtime.
        Both backtracking and multiple results often manifest within \VC's
        choice operator, which tempt its removal altogether. However, we are drawn
        to harnessing the expressive potential of this operator, particularly
        when paired with \VC's 'one' keyword. When employed with choice as a
        condition, 'one' elegantly signifies 'proceed if any branch of the
        choice succeeds. 
        
        To this end, we imagine a new language, \VMinus, in which choice is
        permitted with several modifications:

        \begin{enumerate}
        \item Choice may only appear as a condition or 'guard', not as a result
        or the right-hand side of a binding.
        \item If any branch of the choice succeeds, the choice succeeds,
        producing any bindings found in that branch. The program examines the
        branches in a left-to-right order.
        \item The existential $\exists$ may not appear under choice.
        \end{enumerate}

        We introduce one more crucial modification to the \VC\ runtime: a name in
        \VMinus is an \it{expression} rather than a \it{value}. This alteration,
        coupled with our adjustments to choice, eradicates backtracking. Our
        rationale behind this is straightforward: if an expression returns a
        name, and subsequently, the program imposes a new constraint on that
        name, it may necessitate the reevaluation of the earlier expression—- a
        scenario we strive to avoid. 
        


        
        
        


        Indeed, having stripped out the functional logic programming elements of
        \VC, only the decision-making bits are left over. To wrap these, we add
        a classic decision-making construct: guarded commands\cite{dijkstra} The
        result is \VMinus. \VMinus's concrete syntax is defined in
        Figure~\ref{fig:vmsyntax}. 

        

    \bf{Programming in \VMinus}
        
    Even with multiple modifications, \VMinus\ still allows for many of the same
    pleasing computations as full Verse. Programmers in \VMinus can\dots
        \begin{enumerate}
            \item Introduce a set of equations, to be solved in a
            nondeterministic order 
            \item Guard expressions with those equations 
            \item Flexibly express “proceed when any of these operations
            succeeds” with the new semantics of choice.  
        \end{enumerate}

    Figure~\ref{fig:vminusfuncs} provides an example of how a programmer can
    utilize \VMinus\ to solve the previous problems: 
    
    
    \begin{figure}[ht] 
      \begin{minipage}[h]{0.54\linewidth}
        \vmlst 
        \begin{lstlisting}[numbers=none, basicstyle=\tiny, xleftmargin=.2em,
          showstringspaces=false,
          frame=single]
val exclaimTall = \sh.
  if sh = Square s; (> s) 100 -> 
            Wow! That's A Tall Square!  
    [] E w h. sh = Triangle w h; (> s) 100 ->
            Goodness! Towering Triangle!
    [] E b1 b2 h. sh = Trapezoid b1 b2 h; 
     (> s) 100 -> Zounds! Tremendous Trapezoid!
    [] ->  Your Shape Is Small
  fi 
  \end{lstlisting}
      \Description{exclaimTall in VMinus}
      \subcaption{\tt{exclaimTall} in \VMinus}
          \label{fig:vmexclaimtall} 
      \end{minipage}%
      \begin{minipage}[h]{0.5\linewidth}
        \vmlst 
        \begin{lstlisting}[numbers=none, basicstyle=\tiny, xleftmargin=2em,
                      frame=single]
  val tripleLookup = \ rho. \x.
    if E x w y z v1 v2 v3. 
      v1 = (lookup rho) x; v1 = Some w; 
      v2 = (lookup rho) w; v2 = Some y; 
      v3 = (lookup rho) y; v3 = Some z; -> 
            z 
      [] -> handleFailure x
    fi 
   \end{lstlisting}
        \Description{tripleLookup in VMinus}
        \subcaption{\tt{tripleLookup} in \VMinus}
            \label{fig:vmtriplelookup} 
        \vspace{4ex}
      \end{minipage} 
      \begin{minipage}[h]{\linewidth}
        \vmlst 
        \begin{lstlisting}[numbers=none, basicstyle=\tiny, xleftmargin=9em,
          showstringspaces=false,
          frame=single]
val game_of_token = \t. 
  if E f. (t = BattlePass f  | t = ChugJug f | t = TomatoTown f) -> 
        P (Fortnite f)
  [] E f. (t = HuntersMark f | t = SawCleaver f)                 -> 
        P (Bloodborne ((* 2) f))
  [] E f. (t = MoghLordOfBlood f | t = PreatorRykard f)          -> 
        P (EldenRing ((* 3) f))
  [] -> P (Irrelevant 0)
  fi 
\end{lstlisting}
      \Description{game_of_token in VMinus}
      \subcaption{\tt{game\_of\_token} in \VMinus}
          \label{fig:vmgot}
      \vspace{4ex}
      \end{minipage}%
      \caption{The functions in \VMinus\ have a desirably concise
      implementation, as before.}
  \label{fig:vminusfuncs}
    \end{figure}   

    \VMinus\ looks satisfyingly similar to both \PPlus\ and \VC. 

    \subsection{\VMinus\ and \PPlus, side by side}
    \rab{include section?}








\subsection{Formal Semantics of ${\VMinus}$:}
\rab{How much semantics should I present here?}
    
    \subsection{Expressions}
    
    \newcommand\GNoTree{\vmrung \rightsquigarrow \uppsidown} 
    
    An expression in core Verse evaluates to produce possibly-empty sequence of
    values. In {\VMinus}, expressions evaluate to a single value. 

    An expression evaluates to produce a \bf{result}. A result is either a
    value $\v$ or \fail. 
    
    \[\result\; \rm{::=}\; \rhohat\; \vert \; \fail \]
    
    \showvjudgement{Eval}{\vmeval}
    \showvjudgement{Process-Guards}{\vmgs}
\subsection{Guards}

In {\VMinus}, we solve guards similarly to how the authors of the original Verse
paper solve equations: we pick one, attempt to solve for it, and move on. 

In our semantics, we do this with a \it{list} of guards: we pick a guard,
attempt to solve it to refine our environment or fail, and repeat. We only pick
a guard out of the context \context\ that we know we can solve. If we can't pick
a guard and there are guards left over, the semantics gets stuck.

Given an environment $\rho$ from names to values {\v}s plus $\bot$, solving a
guard \g\ will either refine the environment ($\Rhoprime$) or be \bf{rejected}.
We use the metavariable $\dagger$ to represent rejection, and a guard producing
$\dagger$ will cause the top-level list of guards to also produce $\dagger$. 

    \showvjudgement{Guard-Refine}{\grefine[result=\rho']}
    \showvjudgement{Guard-reject}\gfail


\rab{I put vminus semantics here, but they're quite long. How much to show?}

\subsection{\VMinus compiles to a decision tree in the ~\D language}
\label{d}

Like pattern matching, \VMinus~compiles to a decision tree whose height is 
linear with respect to code size. To formalize decision trees, we use 
\D, a language which generalizes the trees found in~\citet{maranget}. 

The spec of decision trees and match compilation algorithm is presented below. 





  
      









\rab{do we want D's concrete syntax?}
  % \begin{figure}
  %   \begin{center}
  %   \dcsyntax
  %   \end{center}
  %   \Description{The concrete syntax of D.}
  %   \caption{\D: Concrete syntax}
  %   \label{fig:dsyntax}
  %   \end{figure}



  In~\D, as in Maranget's trees, the~\it{test} node extracts all names from a
  value constructor at once for use in subtrees. The compiler is responsible for
  introducing the fresh names used in~\it{test}. The compiler alpha-renames all
  necessary terms before it translates an~\iffibf to a decision tree to ensure
  all names are unique. 

  In~\D, like in~\VMinus, expressions can~\it{fail}, meaning some of~\D's
  syntactic forms like~\it{try-let} and~\it{cmp} have an extra branch that is
  executed if the examined expression fails. 

\rab{d semantics- show or omit?}


    \subsection{The~\DTran\ algorithm:~\VMinus $\rightarrow$~\D}



       




    When presented with an~\iffibf,~\DTran\~invokes~\Compile with context~\ctx.
    \Compile~first desugars the~\iffibf by expanding choice to multiple~\iffibf
    branches with a desugaring function~\ITran: 
    
    \itran{if\;~\dots~\dbar\; gs_{1};\;~\choiceg{gs_{2}}{gs_{3}};\; gs_{4}~\rightarrow e\;~\dbar~\dots~\;fi}
    
    ==
    
    ${if\;~\dots~\dbar\; gs_{1};\; gs_{2}~\rightarrow e~\;\dbar\; gs_{3};\; gs_{4}~\rightarrow e~\;\dbar~\dots~\;fi}$

    With the desugared~\iffibf,~\Compile\ then repeatedly chooses a guarded
    expression $G$ and applies one of the compilation rules below. The rules are
    applied in a nondeterministic order. 

    \Compile terminates when it inserts a final~\it{match} node (with rule
    \textsc{Match}).  A~\it{match} node is inserted for a right-hand side expression~\expr when the list of
    guards preceding~\expr is empty or a list of assignments from names to
    unbound names. Termination of~\DTran\~is guaranteed because each recursive
    call passes a list of guarded expressions in which the number of guards is
    strictly smaller, so eventually the algorithm reaches a state in which the
    first unmatched branch is all trivially-satisfied guards.


    If~\Compile\ cannot choose a $g$ of one of the valid forms, it halts with an
    error. This can happen when no $g$ is currently solvable in the context, as
    determined by the same algorithm that~\VMinus uses to pick a guard to solve,
    or when the program would be forced to unify incompatible values, such as
    any value with a closure. 

    \rab{I have these rules with english descriptions, though that increases length}.

    
\showvjudgement{Compile}{\compilebig}

      \subsection{Reduction Strategies}
      \rab{really interesting section, want to include}
      % In our implementation, we apply the~\textsc{Test} rule before other rules,
      % having found in our experiments that this heuristic leads to the smallest
      % trees. A~%
      % risk of inserting a \it{test} node before a \it{let-unless} node is that
      % if there are many shared names across branches, a \it{test} node will
      % introduce those names in \it{let-unless} nodes in each subtree. It is
      % possible to insert \it{let-unless} nodes before \it{test} nodes, in which
      % case those names will be bound only once.  The risk of the \it{let}-first
      % strategy is that if there are many names used in only one subtree of a
      % \it{test}, these names will be introduced to \it{all} branches.  The extra
      % bindings may harm performance by bloating an environment in an interpreter
      % or thrashing the icache if the tree further is compiled to machine code. 

      % In the implementation,~\DTran\~first introduces all the names under the
      % all existential $\exists$'s to a context which determines if a name is
      % \it{known} or~\it{unknown}, applying the~\textsc{Exists} rule is applied
      % all at the start. At the start of compilation, each name introduced by
      % $\exists$ is~\it{unknown} in a context~\ctx. Since all names in the
      % program are unique at this stage, there are no clashes. 

      \subsection{Full Big-step rules for \Compile, with no descriptions}

        The judgement form for compilation is \textsc{Compile}: 

        \showvjudgement{Compile}{\compilebig}

The rules are nondeterministic: the structure of the final result~\expr\ depends
on the order in which rules are applied by the compiler.  

\begingroup
\scriptsize
\rawcompiler
\endgroup


    \section{Implementations}

    The implementations of \VMinus, \PPlus, and~\D are complete, from parsers to
    evaluation to unparsers. In the same repository lives the complete
    implementation of~\DTran, which translates from~\VMinus to~\D. With the
    implementations, we~include test cases for evaluation of~\VMinus, evaluation
    of~\D, and the translation between the two. In~each of these test cases, the
    translation preserves semantics. 


    \section{Related and Future Work}

    This paper builds on Augustsson et al.'s Verse Calculus~\citep{verse} and
    decision trees~\citep{maranget}. Augustsson et al. give the formal rewrite
    semantics for the Verse Calculus; Maranget gives an elegant formalism of
    decision trees and a translation algorithm from patterns to decision trees.
    We~attempted to imitate the behavior of the rewrite semantics of~\VC in the
    big-step semantics of~\VMinus by manually rewriting terms in~\VC and
    creating rules that would imitate the ultimate result of term-rewriting.
    
    \rab{"big-step-looking interpreters are the way we like to write them?"
    what's a good way to say this point? } 
    
    We~chose a big-step semantics because
    the big-step evaulation rules closely mirrors the structure of our
    interpreters and compiler, which helped us reason about the code. 

    \rab{axe below, because too related to Verse?}

    Using a rewrite semantics instead would more closely relate~\VMinus
    and~\VC, and is a likely future project. 
    
    Maranget's formalism was the foundation off of which we~built~\D. 
    
    Extensions to pattern matching, and how they appeal to language designers,
    find an excellent example in~\citet{guardproposal}. The authors describe
    pattern guards and transformational patterns (another extension to pattern
    matching), both of which allow a Haskell programmer to write more concise
    code using pattern matching. Or-patterns are documented in the OCaml
    Language Reference Manual~\citep{ocaml}.
    
    \citet{augustsson1985compiling}~ gives a foundation in compiling pattern
    matching. \citet{scottramsey} have a crisp example of a match-compilation
    algorithm (pattern matching to decision trees). Scott and Ramsey's algorithm
    structurally inspired mine, and studying the source code from the paper
    aided our implementation. 

    For future work, our top priority is to prove that \DTran~preserves
    semantics. Secondarily, we~plan to prove that \VMinus, like pattern
    matching, is deterministic. 

    Finally, we~would like to explore
    exhaustiveness analysis of~\VMinus.
    Exhaustiveness analysis
    can warn programmers of a missing or extraneous alternative in a~\it{case} expression.
    Owing to its significantly more flexible structure,
    however,~\iffibf may prove trickier to analyze.
    
    
    
    
    
    
    
    
    
    
    
    
    
    
    
    
    
        
    


    
    
    

    
    
        







        

            
            
            





        
        
        

        

    \section{Conclusion}

    We~have introduced \VMinus, a~language that makes decisions using equations.
    By~example, we have shown that programs written in~\VMinus can have the same
    desirable properties as equivalent programs written using pattern matching.
    And~to~show that equations can be compiled to efficient decision trees, we
    have introduced~\D and~\DTran. In doing so, we~have demonstrated that
    programming with equations is a promising alternative to pattern matching.
    
    



\end{document}
